\documentclass[a4paper,12pt]{article}

\usepackage[T2A]{fontenc}
\usepackage[utf8]{inputenc}
\usepackage[russian]{babel}
\usepackage{amsmath,amssymb,mathtools}

\usepackage{helvet}
% PTSans обеспечивает корректный кириллический sans-serif в pdfLaTeX.
\usepackage{PTSans}
\renewcommand{\familydefault}{\sfdefault}

\usepackage{icomma}
\usepackage[
  a4paper,
  left=20mm,
  right=20mm,
  top=20mm,
  bottom=22mm
]{geometry}

\usepackage{microtype}
\usepackage{tikz}
\usetikzlibrary{calc,arrows.meta,positioning,shapes.geometric}

\usepackage{tabularx,booktabs}
\usepackage{enumitem}
\usepackage{hyperref}
\usepackage{parskip}
\usepackage{fancyhdr}
\usepackage{setspace}
\usepackage{xcolor}

% ------------------------------------------------------------
% Палитра
% ------------------------------------------------------------

\definecolor{accent}{HTML}{008080}
\definecolor{darkaccent}{HTML}{004D4D}
\definecolor{textgray}{HTML}{333333}
\definecolor{softgray}{HTML}{F3F5F5}

% ------------------------------------------------------------
% Общая типографика
% ------------------------------------------------------------

\hypersetup{hidelinks}

\onehalfspacing
\setlength{\parindent}{0pt}
\setlength{\headheight}{25pt}
\setlength{\emergencystretch}{2em}

\renewcommand{\arraystretch}{1.35}

\setlist[itemize]{
  leftmargin=6mm,
  itemsep=2pt,
  topsep=3pt
}

\setlist[enumerate]{
  leftmargin=8mm,
  itemsep=4pt,
  topsep=4pt
}

% ------------------------------------------------------------
% Колонтитулы
% ------------------------------------------------------------

\pagestyle{fancy}
\fancyhf{}

\lhead{\small Анастасия Павлова}
\rhead{\small ЕГЭ: прикладные задачи со степенями}
\cfoot{\small \thepage}

\renewcommand{\headrulewidth}{0.4pt}
\renewcommand{\headrule}{%
  \hbox to\headwidth{%
    \color{accent}%
    \leaders\hrule height \headrulewidth\hfill
  }%
}

% ------------------------------------------------------------
% TikZ: единые стили блок-схем
% ------------------------------------------------------------

\tikzset{
  flowstart/.style={
    ellipse,
    draw=darkaccent,
    line width=0.9pt,
    minimum width=42mm,
    minimum height=11mm,
    align=center,
    font=\small
  },
  flowproc/.style={
    rectangle,
    rounded corners=2mm,
    draw=accent,
    line width=0.9pt,
    minimum height=12mm,
    text width=112mm,
    align=center,
    font=\small,
    inner sep=4mm
  },
  flowarrow/.style={
    -{Stealth[length=3mm,width=2mm]},
    line width=0.9pt,
    draw=darkaccent
  }
}

% ------------------------------------------------------------
% Служебные команды
% ------------------------------------------------------------

\newcommand{\sect}[1]{%
  \section*{\color{darkaccent}#1}%
}

\newcommand{\subsect}[1]{%
  \subsection*{\color{accent}#1}%
}

\newcommand{\key}[1]{%
  \textbf{\color{darkaccent}#1}%
}

% ============================================================
% Документ
% ============================================================

\begin{document}

% ============================================================
% ТИТУЛЬНЫЙ ЛИСТ
% ============================================================

\begin{titlepage}

\thispagestyle{empty}

\centering

\vspace*{24mm}

{\Large
\color{darkaccent}
\textbf{Чек-лист}
}

\vspace{8mm}

{\huge
\textbf{Прикладные задачи со степенями}
}

\vspace{4mm}

{\Large
\textbf{модели, граничные значения и параметры}
}

\vspace{16mm}

{\large
Подготовка к ЕГЭ по профильной математике
}

\vspace{12mm}

{\large
Анастасия Павлова
}

\vspace{18mm}

{\large
\today
}

\vspace{18mm}

{\normalsize
Лёвин Артём Александрович\\
эксклюзивно для Анастасии Павловой
}

\vfill

\begin{tikzpicture}[remember picture,overlay]

\draw[
  accent,
  line width=1.15pt
]
  ([xshift=18mm,yshift=18mm]current page.south west)
  ..
  controls
  ([xshift=55mm,yshift=31mm]current page.south west)
  and
  ([xshift=-58mm,yshift=8mm]current page.south east)
  ..
  ([xshift=-18mm,yshift=20mm]current page.south east);

\end{tikzpicture}

\end{titlepage}

% ============================================================
% ОСНОВНОЙ ЧЕК-ЛИСТ
% ============================================================

\sect{1. Что требуется уметь после занятия}

\begin{itemize}

  \item распознавать прикладную модель и выделять математически существенные данные;

  \item корректно работать с условиями
  \textit{не меньше},
  \textit{не больше},
  \textit{наибольший},
  \textit{наименьший};

  \item переводить числа из десятичной записи в удобный степенной вид;

  \item выражать неизвестную величину из формулы;

  \item корректно работать с дробными степенями;

  \item решать показательные уравнения приведением к одному основанию;

  \item понимать запись вида $M(t)$ как значение функции, зависящей от $t$;

  \item переводить фразы
  «уменьшилось в $q$ раз»
  и
  «увеличилось в $r$ раз»
  в математические соотношения;

  \item сравнивать два состояния одной системы при условии
  \[
  pV^a=\mathrm{const}.
  \]

\end{itemize}

% ------------------------------------------------------------

\subsect{Базовые ограничения}

В рассматриваемых моделях давление, объём, масса и характерные
параметры положительны.

Для модели
\[
pV^k=C
\]
используем условия
\[
p>0,
\qquad
V>0,
\qquad
C>0,
\qquad
k>0.
\]

Для модели распада
\[
M(t)=M_0b^{-t/T}
\]
предполагаем
\[
M_0>0,
\qquad
b>1,
\qquad
T>0,
\qquad
t\ge0.
\]

Единицы измерения не влияют на алгебраические преобразования.
Формат окончательного ответа должен соответствовать условию задачи.

% ============================================================
% МОДЕЛЬ pV^k = C
% ============================================================

\sect{2. Степенная модель \texorpdfstring{$pV^k=C$}{pV^k=C}}

\subsect{Главная идея}

Пусть
\[
pV^k=C.
\]

Выражаем степенную часть:

\[
V^k=\frac{C}{p}.
\]

Следовательно,

\[
V=
\left(
\frac{C}{p}
\right)^{1/k}.
\]

При
\[
C>0,
\qquad
k>0
\]
величина $V$ уменьшается при увеличении $p$.

Поэтому

\[
p\ge p_0
\quad\Longrightarrow\quad
V\le
\left(
\frac{C}{p_0}
\right)^{1/k}.
\]

Следовательно, \key{наибольший допустимый объём}
достигается при граничном значении

\[
p=p_0.
\]

Именно монотонность позволяет перейти от условия с неравенством
к граничному равенству.

% ------------------------------------------------------------

\subsect{Работа с десятичными коэффициентами}

Перед степенными преобразованиями удобно избавиться от
десятичной дроби в коэффициенте.

Например,

\[
6{,}25\cdot10^5
=
625\cdot10^3.
\]

Если запятую сдвинули вправо на $r$ разрядов,
показатель степени десяти уменьшается на $r$:

\[
a\cdot10^n
=
(10^ra)\cdot10^{n-r}.
\]

Например,

\[
3{,}125\cdot10^7
=
3125\cdot10^4.
\]

% ------------------------------------------------------------

\subsect{Как убрать дробную степень}

Пусть

\[
V^{5/3}=A,
\qquad
A>0.
\]

Требуется получить первую степень $V$.

Используем показатель, обратный числу $5/3$:

\[
\frac{3}{5}.
\]

Возводим обе части в степень $3/5$:

\[
\left(
V^{5/3}
\right)^{3/5}
=
A^{3/5}.
\]

Слева

\[
V^{(5/3)(3/5)}
=
V^1
=
V.
\]

Итак,

\[
\boxed{
V=A^{3/5}
}.
\]

Ключевая проверка:

\[
\frac53\cdot\frac35=1.
\]

% ------------------------------------------------------------

\subsect{Пример}

Пусть

\[
pV^{5/3}
=
2{,}43\cdot10^6,
\qquad
p\ge10^4.
\]

Требуется найти максимальный объём.

Поскольку при увеличении $p$ объём уменьшается,
для максимума используем

\[
p=10^4.
\]

Подставляем:

\[
10^4V^{5/3}
=
2{,}43\cdot10^6.
\]

Отсюда

\[
V^{5/3}
=
\frac{2{,}43\cdot10^6}{10^4}
=
243.
\]

Представляем число степенью:

\[
243=3^5.
\]

Получаем

\[
V^{5/3}=3^5.
\]

Возводим обе части в степень $3/5$:

\[
V
=
(3^5)^{3/5}
=
3^3
=
27.
\]

\[
\boxed{V_{\max}=27}
\]

% ------------------------------------------------------------

\subsect{Типичные ошибки}

\begin{itemize}

  \item заменять исходное неравенство равенством без проверки монотонности;

  \item при $V^{5/3}=A$ возводить части уравнения в степень $5/3$;

  \item забывать, что необходима обратная степень:
  \[
  \frac53\longrightarrow\frac35;
  \]

  \item оставлять составное число в неудобном виде, хотя
  \[
  243=3^5,
  \qquad
  3125=5^5,
  \qquad
  32=2^5;
  \]

  \item терять скобки в выражениях вида
  \[
  \left(\frac AB\right)^r;
  \]

  \item выполнять вычисления с десятичными дробями там,
  где переход к степенному виду существенно упрощает решение.

\end{itemize}

% ============================================================
% FLOWCHART 1
% ============================================================

\clearpage
\thispagestyle{empty}

\begin{center}

{\Large
\color{darkaccent}
\textbf{
Алгоритм: максимальный объём в модели $pV^k=C$
}
}

\end{center}

\vspace{8mm}

\begin{center}

\begin{tikzpicture}[node distance=7mm]

\node[flowstart] (s)
{Начало};

\node[flowproc,below=of s] (a)
{
Запишите модель
$pV^k=C$
и ограничения
$C>0$, $k>0$, $p>0$, $V>0$.
};

\node[flowproc,below=of a] (b)
{
Выразите объём:
\[
V=
\left(
\frac Cp
\right)^{1/k}.
\]
};

\node[flowproc,below=of b] (c)
{
Используйте монотонность:
при $k>0$ объём $V$ уменьшается при росте $p$.
};

\node[flowproc,below=of c] (d)
{
Если $p\ge p_0$ и требуется $V_{\max}$,
подставьте граничное значение
$p=p_0$.
};

\node[flowproc,below=of d] (e)
{
Решите
\[
V^k=\frac{C}{p_0}.
\]
Приведите числа к удобному степенному виду.
};

\node[flowproc,below=of e] (f)
{
Возведите обе части в степень $1/k$
и получите значение $V$.
};

\node[flowproc,below=of f] (g)
{
Проверьте $V>0$, исходное ограничение
и требуемый формат ответа.
};

\node[flowstart,below=of g] (t)
{Ответ};

\draw[flowarrow] (s)--(a);
\draw[flowarrow] (a)--(b);
\draw[flowarrow] (b)--(c);
\draw[flowarrow] (c)--(d);
\draw[flowarrow] (d)--(e);
\draw[flowarrow] (e)--(f);
\draw[flowarrow] (f)--(g);
\draw[flowarrow] (g)--(t);

\end{tikzpicture}

\end{center}

\clearpage

% ============================================================
% ЭКСПОНЕНЦИАЛЬНАЯ МОДЕЛЬ
% ============================================================

\sect{3. Показательная модель распада}

\subsect{Как читать запись \texorpdfstring{$M(t)$}{M(t)}}

Запись

\[
M(t)=M_0\,2^{-t/T}
\]

показывает зависимость массы $M$ от времени $t$.

Скобки в записи $M(t)$ обозначают аргумент функции.
Произведения $M\cdot t$ здесь нет.

Обозначения:

\begin{itemize}

  \item $M_0$ --- начальная масса;

  \item $T$ --- положительная постоянная модели;

  \item $t$ --- время;

  \item $M(t)$ --- масса в момент времени $t$;

  \item множитель
  \[
  2^{-t/T}
  \]
  задаёт убывание величины.

\end{itemize}

% ------------------------------------------------------------

\subsect{Граничное значение времени}

При $T>0$ функция

\[
M(t)=M_0\,2^{-t/T}
\]

убывает.

Поэтому условие

\[
M(t)\ge m_*
\]

задаёт промежуток допустимых значений времени.

Если требуется определить,
в течение какого времени масса остаётся не меньше $m_*$,
последний допустимый момент находится из граничного равенства

\[
M(t)=m_*.
\]

После нахождения границы полезно сформулировать итоговый интервал:

\[
0\le t\le t_{\max}.
\]

% ------------------------------------------------------------

\subsect{Пример}

Пусть

\[
M(t)=40\cdot2^{-t/10},
\qquad
M(t)\ge5.
\]

На границе допустимого промежутка

\[
40\cdot2^{-t/10}=5.
\]

Делим обе части на $40$:

\[
2^{-t/10}
=
\frac5{40}
=
\frac18.
\]

Представляем дробь степенью двойки:

\[
\frac18=2^{-3}.
\]

Получаем

\[
2^{-t/10}
=
2^{-3}.
\]

Основания одинаковы, поэтому

\[
-\frac{t}{10}
=
-3.
\]

Умножаем обе части на $-10$:

\[
t=30.
\]

Следовательно,

\[
0\le t\le30.
\]

Последний допустимый момент:

\[
\boxed{t=30}.
\]

% ------------------------------------------------------------

\subsect{Полезные формулы для выражения неизвестного множителя}

Если

\[
ab=c
\]

и $a\ne0$, то

\[
b=\frac ca.
\]

Аналогично, при $b\ne0$:

\[
a=\frac cb.
\]

Если

\[
\frac ab=c,
\qquad
b\ne0,
\]

то

\[
a=bc.
\]

Если дополнительно $c\ne0$, то

\[
b=\frac ac.
\]

Простейший шаблон:

\[
40X=2{,}5
\]

даёт

\[
X=\frac{2{,}5}{40}.
\]

Порядок деления здесь принципиален.

% ------------------------------------------------------------

\subsect{Показательное уравнение: основной принцип}

Если получено уравнение

\[
a^{f(t)}=a^{g(t)},
\qquad
a>0,
\qquad
a\ne1,
\]

то

\[
f(t)=g(t).
\]

Поэтому при работе со степенями полезно сразу искать возможность
представить составные числа и дроби через одно основание:

\[
8=2^3,
\qquad
16=2^4,
\qquad
\frac18=2^{-3},
\qquad
\frac1{16}=2^{-4}.
\]

% ------------------------------------------------------------

\subsect{Типичные ошибки}

\begin{itemize}

  \item менять местами делимое и делитель;

  \item добавлять знак «минус» при переносе множителя;

  \item записывать
  \[
  40X=2{,}5
  \]
  как
  \[
  X=\frac{40}{2{,}5};
  \]

  \item забывать правило
  \[
  \frac1{2^n}=2^{-n};
  \]

  \item получать отрицательное время и не выполнять проверку;

  \item путать постоянную $T$ с искомой переменной $t$;

  \item воспринимать $M(t)$ как произведение двух величин.

\end{itemize}

% ============================================================
% FLOWCHART 2
% ============================================================

\clearpage
\thispagestyle{empty}

\begin{center}

{\Large
\color{darkaccent}
\textbf{
Алгоритм: последнее допустимое время в модели распада
}
}

\end{center}

\vspace{8mm}

\begin{center}

\begin{tikzpicture}[node distance=7mm]

\node[flowstart] (s)
{Начало};

\node[flowproc,below=of s] (a)
{
Запишите
$M(t)=M_0b^{-t/T}$,
где
$M_0>0$, $b>1$, $T>0$, $t\ge0$.
};

\node[flowproc,below=of a] (b)
{
Переведите условие
«масса не меньше $m_*$»
в неравенство
$M(t)\ge m_*$.
};

\node[flowproc,below=of b] (c)
{
Учтите убывание функции:
последнее допустимое время достигается на границе
$M(t)=m_*$.
};

\node[flowproc,below=of c] (d)
{
Разделите обе части на $M_0$:
\[
b^{-t/T}
=
\frac{m_*}{M_0}.
\]
};

\node[flowproc,below=of d] (e)
{
Представьте правую часть степенью основания $b$.
};

\node[flowproc,below=of e] (f)
{
Приравняйте показатели
и решите полученное линейное уравнение относительно $t$.
};

\node[flowproc,below=of f] (g)
{
Проверьте $t\ge0$
и подстановкой убедитесь,
что граничное значение найдено верно.
};

\node[flowstart,below=of g] (t)
{Ответ};

\draw[flowarrow] (s)--(a);
\draw[flowarrow] (a)--(b);
\draw[flowarrow] (b)--(c);
\draw[flowarrow] (c)--(d);
\draw[flowarrow] (d)--(e);
\draw[flowarrow] (e)--(f);
\draw[flowarrow] (f)--(g);
\draw[flowarrow] (g)--(t);

\end{tikzpicture}

\end{center}

\clearpage

% ============================================================
% ПАРАМЕТР a
% ============================================================

\sect{4. Сравнение двух состояний:
\texorpdfstring{$pV^a=\mathrm{const}$}{pV^a=const}}

\subsect{Переход к индексам}

Если одна система подчиняется закону

\[
pV^a=\mathrm{const},
\]

то для двух её состояний справедливо

\[
p_1V_1^a
=
p_2V_2^a.
\]

Индексы обозначают состояния:

\[
(p_1,V_1)
\]

--- исходное состояние,

\[
(p_2,V_2)
\]

--- новое состояние.

Параметр $a$ остаётся одним и тем же.

% ------------------------------------------------------------

\subsect{Перевод словесных изменений}

Если объём уменьшился в $q$ раз, то

\[
V_2
=
\frac{V_1}{q}.
\]

Эквивалентная запись:

\[
V_1=qV_2.
\]

Вторая форма часто удобнее при подстановке.

Если давление увеличилось в $r$ раз, то

\[
p_2=rp_1.
\]

Например:

\[
\text{объём уменьшился в 2 раза}
\quad\Longrightarrow\quad
V_1=2V_2,
\]

\[
\text{давление увеличилось в 4 раза}
\quad\Longrightarrow\quad
p_2=4p_1.
\]

% ------------------------------------------------------------

\subsect{Пример}

Пусть

\[
pV^a=\mathrm{const}.
\]

Объём уменьшился в $2$ раза,
а давление увеличилось в $4$ раза.

Записываем:

\[
V_1=2V_2,
\qquad
p_2=4p_1.
\]

Для двух состояний:

\[
p_1V_1^a
=
p_2V_2^a.
\]

Подставляем:

\[
p_1(2V_2)^a
=
(4p_1)V_2^a.
\]

Используем степень произведения:

\[
(2V_2)^a
=
2^aV_2^a.
\]

Получаем

\[
p_1\,2^aV_2^a
=
4p_1V_2^a.
\]

Поскольку

\[
p_1>0,
\qquad
V_2>0,
\]

общие множители можно сократить:

\[
2^a=4.
\]

Представляем $4$ степенью двойки:

\[
4=2^2.
\]

Следовательно,

\[
2^a=2^2,
\]

откуда

\[
\boxed{a=2}.
\]

% ------------------------------------------------------------

\subsect{Правило подстановки}

Составное подставляемое выражение следует брать в скобки.

При

\[
V_1=2V_2
\]

в выражении $V_1^a$ получаем

\[
V_1^a
\longrightarrow
(2V_2)^a.
\]

Затем используется правило степени произведения:

\[
(xy)^r=x^ry^r.
\]

Например,

\[
(3V_2)^a
=
3^aV_2^a.
\]

Скобки особенно важны, когда подставляемая сумма,
разность или произведение возводится в степень.

% ------------------------------------------------------------

\subsect{Универсальная модель изменения в несколько раз}

Пусть объём уменьшился в $q$ раз,
а давление увеличилось в $r$ раз.

Тогда

\[
V_1=qV_2,
\qquad
p_2=rp_1.
\]

Подстановка в

\[
p_1V_1^a=p_2V_2^a
\]

даёт

\[
p_1(qV_2)^a
=
rp_1V_2^a.
\]

После раскрытия степени:

\[
p_1q^aV_2^a
=
rp_1V_2^a.
\]

Сокращаем положительные общие множители:

\[
q^a=r.
\]

Таким образом, задача сводится к показательному уравнению

\[
\boxed{q^a=r}.
\]

Если $q$ и $r$ можно представить степенями одного основания,
получаем особенно простой расчёт.

Например,

\[
4^a=16.
\]

Поскольку

\[
4=2^2,
\qquad
16=2^4,
\]

имеем

\[
(2^2)^a=2^4,
\]

\[
2^{2a}=2^4,
\]

\[
2a=4,
\]

\[
a=2.
\]

% ------------------------------------------------------------

\subsect{Типичные ошибки}

\begin{itemize}

  \item путать исходное и новое значение;

  \item неверно переводить уменьшение в $q$ раз:
  \[
  V_2=\frac{V_1}{q};
  \]

  \item записывать увеличение давления в $r$ раз как деление на $r$;

  \item забывать скобки при подстановке;

  \item индексировать параметр $a$ отдельно для каждого состояния;

  \item сокращать множители без проверки их ненулевого значения;

  \item вычислять сложные степени напрямую, хотя выражение можно
  привести к одному основанию.

\end{itemize}

% ============================================================
% FLOWCHART 3
% ============================================================

\clearpage
\thispagestyle{empty}

\begin{center}

{\Large
\color{darkaccent}
\textbf{
Алгоритм: нахождение параметра $a$
}
}

\end{center}

\vspace{8mm}

\begin{center}

\begin{tikzpicture}[node distance=7mm]

\node[flowstart] (s)
{Начало};

\node[flowproc,below=of s] (a)
{
Запишите закон $pV^a=C$ для двух состояний:
\[
p_1V_1^a=p_2V_2^a.
\]
};

\node[flowproc,below=of a] (b)
{
Переведите изменения в формулы:
при уменьшении объёма в $q$ раз
$V_1=qV_2$;
при росте давления в $r$ раз
$p_2=rp_1$.
};

\node[flowproc,below=of b] (c)
{
Подставьте выражения со скобками:
\[
p_1(qV_2)^a=(rp_1)V_2^a.
\]
};

\node[flowproc,below=of c] (d)
{
Раскройте степень произведения:
\[
p_1q^aV_2^a
=
rp_1V_2^a.
\]
};

\node[flowproc,below=of d] (e)
{
Сократите положительные общие множители
$p_1$ и $V_2^a$.
};

\node[flowproc,below=of e] (f)
{
Получите
\[
q^a=r
\]
и приведите обе части к одному основанию.
};

\node[flowproc,below=of f] (g)
{
Приравняйте показатели,
найдите $a$
и проверьте исходное соотношение.
};

\node[flowstart,below=of g] (t)
{Ответ};

\draw[flowarrow] (s)--(a);
\draw[flowarrow] (a)--(b);
\draw[flowarrow] (b)--(c);
\draw[flowarrow] (c)--(d);
\draw[flowarrow] (d)--(e);
\draw[flowarrow] (e)--(f);
\draw[flowarrow] (f)--(g);
\draw[flowarrow] (g)--(t);

\end{tikzpicture}

\end{center}

\clearpage

% ============================================================
% ИТОГОВАЯ ПАМЯТКА
% ============================================================

\sect{5. Итоговая памятка}

\begin{enumerate}

  \item
  \key{Начинайте с рабочей формулы.}

  Выпишите математическую модель и определите,
  какие величины известны.

  \item
  \key{Учитывайте граничные слова.}

  Условия
  «не меньше»,
  «не больше»,
  «наибольший»,
  «наименьший»
  требуют проверки монотонности.

  \item
  \key{Приводите десятичные коэффициенты к удобному виду.}

  Например,

  \[
  6{,}25\cdot10^5
  =
  625\cdot10^3.
  \]

  \item
  \key{Для дробной степени используйте обратный показатель.}

  Если

  \[
  x^{m/n}=A
  \]

  при допустимых положительных значениях, то

  \[
  x=A^{n/m}.
  \]

  \item
  \key{В показательных уравнениях ищите общее основание.}

  Например,

  \[
  \frac1{16}=2^{-4}.
  \]

  \item
  \key{Запись функции содержит аргумент.}

  Выражение $M(t)$ показывает зависимость $M$ от $t$.

  \item
  \key{Уменьшение в $q$ раз связано с делением.}

  \[
  V_{\text{нов}}
  =
  \frac{V_{\text{стар}}}{q}.
  \]

  \item
  \key{Увеличение в $q$ раз связано с умножением.}

  \[
  p_{\text{нов}}
  =
  q\,p_{\text{стар}}.
  \]

  \item
  \key{Составную подстановку заключайте в скобки.}

  \[
  V_1=2V_2
  \quad\Longrightarrow\quad
  V_1^a=(2V_2)^a.
  \]

  \item
  \key{Проверяйте смысл результата.}

  В рассмотренных моделях

  \[
  p>0,
  \qquad
  V>0,
  \qquad
  M>0,
  \qquad
  t\ge0.
  \]

\end{enumerate}

% ============================================================
% ТРЕНИРОВОЧНЫЙ БЛОК
% ============================================================

\clearpage

\sect{Тренировочный блок. Путь Б}

\textbf{Правило работы.}
Решайте по шаблону чек-листа,
обязательно записывая ключевые преобразования.
Задания внутри каждого дня расположены по нарастающей сложности.

% ------------------------------------------------------------

\subsect{11.09.2026}

\begin{enumerate}[label=\arabic*)]

  \item
  Газ подчиняется закону
  \[
  pV^{3/2}=4{,}0\cdot10^7.
  \]
  Найдите наибольший объём $V$, если
  \[
  p\ge5\cdot10^6.
  \]

  \item
  Газ подчиняется закону
  \[
  pV^{5/3}=2{,}43\cdot10^6.
  \]
  Найдите наибольший объём $V$, если
  \[
  p\ge10^4.
  \]

  \item
  Масса вещества меняется по закону
  \[
  M(t)=64\cdot2^{-t/3}.
  \]
  В течение скольких единиц времени масса будет не меньше $8$?

  \item
  Величина меняется по закону
  \[
  A(t)=81\cdot3^{-t/2}.
  \]
  Найдите последнее значение $t\ge0$, при котором
  \[
  A(t)\ge3.
  \]

  \item
  Для газа выполняется
  \[
  pV^a=\mathrm{const}.
  \]
  Объём уменьшили в $2$ раза,
  после чего давление увеличилось в $8$ раз.
  Найдите $a$.

\end{enumerate}

% ------------------------------------------------------------

\subsect{12.09.2026}

\begin{enumerate}[label=\arabic*)]

  \item
  Газ подчиняется закону
  \[
  pV^{4/3}=1{,}62\cdot10^5.
  \]
  Найдите наибольший объём $V$, если
  \[
  p\ge2\cdot10^3.
  \]

  \item
  Газ подчиняется закону
  \[
  pV^{3/2}=2{,}56\cdot10^6.
  \]
  Найдите наибольший объём $V$, если
  \[
  p\ge4\cdot10^4.
  \]

  \item
  Масса вещества меняется по закону
  \[
  M(t)=100\cdot2^{-t/5}.
  \]
  В течение скольких единиц времени масса будет не меньше
  $12{,}5$?

  \item
  Численность модели задаётся формулой
  \[
  N(t)=243\cdot3^{-t/4}.
  \]
  Найдите последнее значение $t\ge0$, при котором
  \[
  N(t)\ge9.
  \]

  \item
  Для газа выполняется
  \[
  pV^a=\mathrm{const}.
  \]
  Объём уменьшили в $4$ раза,
  давление увеличилось в $16$ раз.
  Найдите $a$.

\end{enumerate}

% ------------------------------------------------------------

\subsect{13.09.2026}

\begin{enumerate}[label=\arabic*)]

  \item
  Газ подчиняется закону
  \[
  pV^{2/3}=1{,}44\cdot10^5.
  \]
  Найдите наибольший объём $V$, если
  \[
  p\ge4\cdot10^3.
  \]

  \item
  Газ подчиняется закону
  \[
  pV^{5/2}=6{,}25\cdot10^6.
  \]
  Найдите наибольший объём $V$, если
  \[
  p\ge2\cdot10^3.
  \]

  \item
  Масса вещества меняется по закону
  \[
  M(t)=128\cdot2^{-t/4}.
  \]
  В течение скольких единиц времени масса будет не меньше $4$?

  \item
  Величина меняется по закону
  \[
  B(t)=729\cdot3^{-t/6}.
  \]
  Найдите последнее значение $t\ge0$, при котором
  \[
  B(t)\ge27.
  \]

  \item
  Для газа выполняется
  \[
  pV^a=\mathrm{const}.
  \]
  Объём уменьшили в $3$ раза,
  давление увеличилось в $81$ раз.
  Найдите $a$.

\end{enumerate}

% ============================================================
% ОТВЕТЫ
% ============================================================

\clearpage

\sect{Ответы к тренировочному блоку}

\begin{table}[ht]

\caption{Ответы по дням}

\centering

\begin{tabularx}
{\linewidth}
{@{}l*{5}{>{\centering\arraybackslash}X}@{}}

\toprule

Дата
&
1
&
2
&
3
&
4
&
5
\\

\midrule

11.09.2026
&
$4$
&
$27$
&
$9$
&
$6$
&
$3$
\\

12.09.2026
&
$27$
&
$16$
&
$15$
&
$12$
&
$2$
\\

13.09.2026
&
$216$
&
$25$
&
$20$
&
$18$
&
$4$
\\

\bottomrule

\end{tabularx}

\end{table}

\vspace{8mm}

\textbf{Контроль формата ответа}

\begin{itemize}

  \item
  в заданиях 1--2 каждого дня указано максимальное допустимое
  значение объёма;

  \item
  в заданиях 3--4 указано последнее допустимое
  неотрицательное значение времени;

  \item
  в задании 5 указано значение параметра $a$.

\end{itemize}

\end{document}